%% The "author" command and its associated commands are used to define
%% the authors and their affiliations.
%% Of note is the shared affiliation of the first two authors, and the
%% "authornote" and "authornotemark" commands
%% used to denote shared contribution to the research.
\author{Martin Farkas}
% \authornote{Both authors contributed equally to this research.}
\email{martin.farkas@edu.bme.hu}
\orcid{0000-0003-4260-717X}
\affiliation{%
  \institution{Budapest University of Technology and Economics}
  \city{Budapest}
  \country{Hungary}
}

\author{Kristóf Marussy}
% \authornotemark[1]
\email{marussy@mit.bme.hu}
\orcid{0000-0002-9135-8256}
\affiliation{%
  \institution{Budapest University of Technology and Economics}
  \city{Budapest}
  \country{Hungary}
}

\author{Oszkár Semeráth}
\email{semerath@mit.bme.hu}
\orcid{0000-0002-3592-5105}
\affiliation{%
  \institution{Budapest University of Technology and Economics}
  \city{Budapest}
  \country{Hungary}
}

\author{Zoltán Micskei}
\email{micskeiz@mit.bme.hu}
\orcid{0000-0003-1846-261X}
\affiliation{%
  \institution{Budapest University of Technology and Economics}
  \city{Budapest}
  \country{Hungary}
}

\author{Imre Kocsis}
\email{ikocsis@mit.bme.hu}
\orcid{0000-0002-2792-3572}
\affiliation{%
  \institution{Budapest University of Technology and Economics}
  \city{Budapest}
  \country{Hungary}
}




%%
%% By default, the full list of authors will be used in the page
%% headers. Often, this list is too long, and will overlap
%% other information printed in the page headers. This command allows
%% the author to define a more concise list
%% of authors' names for this purpose.
\renewcommand{\shortauthors}{Farkas et al.}
